\section{Introduction}

Large language models (LLMs) deployed in high-stakes reasoning tasks frequently exploit \textit{reasoning shortcuts}—surface pattern matching, answer memorization, keyword-answer correlations, and premature answer prediction—rather than performing genuine logical inference~\cite{wei2022chain,wang2022self}. This shortcut reliance stems from \textbf{training signal misalignment}: shortcut-encouraging samples efficiently reduce training loss but fail to foster generalizable reasoning. Consequently, models accurate on training distributions fail catastrophically when numbers change, constraints vary, or problem phrasing shifts. This motivates \textit{\textbf{S}hortcut-\textbf{A}ware \textbf{R}easoning \textbf{T}raining (SART)}: identifying samples that promote shortcut updates and modifying gradient dynamics to emphasize genuine reasoning signals. Solving SART is critical for reliable LLM deployment in mathematical reasoning, financial analysis, planning, and scientific discovery—domains where spurious reasoning leads to incorrect decisions or unsafe recommendations.


Two major challenges arise in solving SART: (1) detecting shortcut-promoting samples from gradient behavior, and (2) modifying training dynamics to suppress shortcut learning without degrading performance. First, shortcut samples are often correctly labeled and indistinguishable by conventional quality metrics. They produce strong gradients that efficiently reduce training loss, making loss-based detection unreliable. The detection challenge is: how can we identify training samples whose gradients improve loss but harm reasoning generalization? Second, simply down-weighting or removing detected shortcut samples is insufficient—it destabilizes training and discards useful signal. The suppression challenge is: how can we modify training dynamics to neutralize shortcut gradients while preserving generalizable reasoning signals from the same samples?

Existing methods only partially address SART. Chain-of-Thought (CoT) training~\cite{wei2022chain,kojima2022large} encourages intermediate reasoning traces but does not modify training dynamics: models can generate plausible steps that still rely on shortcuts, and shortcut gradients remain dominant~\cite{ho2022large}. Self-Consistency Decoding~\cite{wang2022self} samples multiple reasoning paths to detect inconsistencies but operates at inference time, leaving shortcut patterns intact in model parameters. Data filtering and curriculum learning~\cite{bengio2009curriculum} attempt to remove noisy samples, but shortcut samples are typically correctly labeled, making them invisible to label-quality filters. RLHF~\cite{ouyang2022training} optimizes reward signals based on answer correctness rather than reasoning validity, allowing shortcuts to persist when they produce correct answers. These approaches cannot jointly address shortcut detection from gradient behavior and dynamic suppression during training—a gradient-centric perspective is required.

\textbf{Our insights: gradient-aware shortcut detection and correction.} We formulate SART as a gradient structure analysis problem. We find that shortcut reasoning samples produce two distinct gradient signatures: (1) \textit{low alignment} with gradients that improve held-out validation performance, indicating the sample's gradient direction does not transfer to generalizable reasoning; and (2) \textit{high answer-token concentration}, where gradient norms are dominated by final answer tokens rather than intermediate reasoning tokens. These signals—gradient alignment with validation gradients and answer-gradient concentration—can be combined into a principled \textbf{ShortcutScore} quantifying the degree to which a sample promotes shortcut learning. Furthermore, harmful shortcut gradient directions can be neutralized via \textit{gradient surgery}: projecting gradients orthogonally onto subspaces that avoid non-transferable directions, forcing the model to learn from genuine reasoning signals without discarding samples entirely.

\textbf{Proposed Approach.} This paper presents the first principled shortcut-aware training framework for LLM reasoning via gradient structure analysis, with two goals: (1) precise detection of shortcut-promoting samples through gradient behavior; (2) effective suppression of shortcut learning while preserving generalizable reasoning signals.
%
For Goal 1, we develop the \textbf{ShortcutScore}, a composite metric combining cosine similarity between per-sample gradients and validation gradients (measuring non-transfer alignment) with the ratio of answer-token to reasoning-token gradient norms (measuring answer-dominant credit assignment). Samples with high ShortcutScore are flagged as shortcut-promoting.
%
For Goal 2, we develop a \textbf{gradient surgery} technique that identifies harmful gradient directions—those misaligned with validation gradients or increasing ShortcutScore—and projects them orthogonally to suppress their influence during parameter updates. This ensures training dynamics prioritize reasoning token signals.

Our framework is model-agnostic and integrates into any LLM training pipeline without architectural changes. Extensive experiments across mathematical reasoning, commonsense reasoning, and constraint-based planning benchmarks demonstrate significant improvements in robustness under distribution shift, number perturbation, and phrasing variation, validating that our approach reduces shortcut reliance and promotes genuine logical inference. Code is available at \url{https://github.com/fuyanjie/short-cut-aware-data-centric-reasoning}.